\bookchapter{الفصل 17: ربط الهاتف بالأجهزة الأخرى}

\chapterintro{الهاتف في يدك هو في الواقع جهاز تحكم عن بُعد خارق ومركز بيانات. فهو لا يجري المكالمات فحسب، بل يربط الحاسوب والتلفاز والساعة ومكبر الصوت والسيارة معاً، بحيث لا تبقى هذه الأجهزة جزراً معزولة. في هذا الفصل سنتحدث عن كيفية جعلها “تشكل فريقاً”.}

\booksection{أولاً: التعاون مع الحاسوب: الهاتف والحاسوب لم يعودا عالمين منفصلين}

لنكن صادقين، كنت أكره قديماً أن ألتقط صورة بالهاتف وأريد تعديلها على الحاسوب، فأضطر للبحث عن كابل بيانات وانتظار النقل. ثم اكتشفت أنك لا تحتاج إلى كابل أصلاً، والربط الجيد أسرع بعشر مرات من الكابل.

\begin{itemize}
  \item \textbf{نظام Apple البيئي (AirDrop + الحافظة العامة)}: إذا كنت تستخدم iPhone مع Mac، فإن AirDrop هي أول ميزة يجب أن تتعلمها. على الهاتف، انقر مشاركة AirDrop اختر Mac الخاص بك، فتنتقل الصورة فوراً، بدون بيانات وبدون Wi-Fi، باستخدام Bluetooth وWi-Fi Direct. وهناك ميزة يسهل تجاهلها اسمها "الحافظة العامة" — تنسخ نصاً على الهاتف، وتلصقه مباشرة على Mac، وكأنه سحر.
  \item \textbf{أندرويد + Windows (Link to Windows)}: لدى Microsoft تطبيق يسمى "Link to Windows"، مدمج في هواتف Samsung، ويمكن لمستخدمي أندرويد الآخرين تنزيله من متجر التطبيقات. بعد الإقران، يمكنك رؤية إشعارات الهاتف والرد على الرسائل النصية وتصفح الصور، بل وإجراء المكالمات مباشرة من الحاسوب. استخدمته لمدة عامين، فلم أعد بحاجة للنظر إلى الهاتف باستمرار أثناء الاجتماعات، تكفي الإشعارات المنبثقة على شاشة الحاسوب.
  \item \textbf{التعاون متعدد الشاشات من العلامات التجارية}: "Multi-Screen Collaboration" من Huawei و"Xiaomi Smart Connect" أكثر جرأة — تعرض شاشة الهاتف مباشرة على الحاسوب، ويمكنك سحب الملفات بالماوس من الهاتف إلى الحاسوب، كما يمكنك تصفح TikTok والرد على WeChat على الحاسوب. النسخ واللصق وسحب الملفات بين الأجهزة ممكنة، والتجربة سلسة جداً.
\end{itemize}

\booksection{ثانياً: العرض على التلفاز: ليشاهد الجميع معاً}

عند زيارة عائلية في عيد رأس السنة، أرادت ابنة عمي عرض فيديو رحلة تخرجها على التلفاز. بعد محاولات طويلة، اكتشفنا أن الأمر بسيط جداً:

\begin{itemize}
  \item \textbf{iPhone التلفاز}: اسحب مركز التحكم من أعلى اليمين، انقر أيقونة المربعين المتداخلين (Screen Mirroring)، ابحث عن اسم تلفازك، وانقر للاتصال. الشرط هو أن يكون التلفاز والهاتف على نفس شبكة Wi-Fi، وأن يدعم التلفاز AirPlay (معظم التلفازات الذكية تدعمه).
  \item \textbf{أندرويد التلفاز}: تختلف التسميات بين العلامات، Huawei تسميه "العرض اللاسلكي"، Xiaomi "الإسقاط"، OPPO/vivo "التفاعل متعدد الشاشات". عموماً تجده في لوحة المفاتيح السريعة، وإن لم تجده فابحث عن "العرض" في الإعدادات.
  \item \textbf{الإسقاط المباشر من تطبيقات الفيديو}: إذا كنت تريد فقط مشاهدة iQiyi أو Bilibili أو YouTube على التلفاز، فالأمر أبسط — ابحث في واجهة التشغيل عن أيقونة صغيرة على شكل مستطيل مع رمز Wi-Fi، انقر عليها، اختر التلفاز، وينتقل الفيديو. حينها يمكنك استخدام الهاتف لأمور أخرى بينما يستمر الفيديو على التلفاز دون تأثير متبادل.
  \item \textbf{ماذا لو فشل العرض}: 90\% من الحالات تكون بسبب عدم الاتصال بنفس شبكة Wi-Fi. و10\% بسبب عدم تفعيل ميزة العرض على التلفاز، ابحث عن مفتاح "التفاعل متعدد الشاشات" أو "AirPlay" في إعدادات التلفاز.
\end{itemize}

\booksection{ثالثاً: مع الساعات الذكية/الأساور: مدير صحة على معصمك}

اشتريت لأمي أساوراً ذكياً، قالت في البداية “أنا لا أتمرن، لماذا أشتري هذا؟”. بعد أسبوعين، اكتشفت أنها تنام خمس ساعات فقط يومياً، فبدأت بتعديل روتين نومها بنفسها. هذه هي القيمة العملية للساعات/الأساور — تحويل بياناتك الصحية غير المرئية إلى أرقام مرئية.

\begin{itemize}
  \item \textbf{فلتر الإشعارات}: عندما يهتز الهاتف في الحقيبة أو الجيب باستمرار، ارفع معصمك للنظرة — رسالة WeChat من المدير أم عرض ترويجي من Pinduoduo، تحكم في ثانية، والرسائل التي لا تستحق إخراج الهاتف دعها تمر.
  \item \textbf{مزامنة البيانات الصحية}: تراقب الساعة على مدار 24 ساعة معدل ضربات القلب، والأكسجين في الدم، والنوم، وعدد الخطوات، وتُجمّع هذه البيانات تلقائياً في تطبيق الصحة على الهاتف. مع الوقت ترى اتجاهاتك — هل مشيت أقل هذا الشهر، هل تحسنت جودة نومك هذا الأسبوع.
  \item \textbf{مزامنة بطاقة المواصلات/بطاقة الدخول}: بطاقات المواصلات وبطاقات الدخول المرتبطة على الهاتف، معظم العلامات تزامنها تلقائياً إلى الساعة. عند ركوب المترو لا تحتاج لإخراج الهاتف، بل قرّب معصمك من بوابة الدخول، وستشعر فعلاً بأن "المستقبل قد وصل".
  \item \textbf{أداة العثور على الهاتف}: معظم الساعات بها ميزة "العثور على الهاتف"، تنقرها فيرنّ الهاتف. أستخدم هذه الميزة بتردد أعلى بكثير مما توقعت — أضع الهاتف في أي مكان ثم أنسى أين هو.
\end{itemize}

\booksection{رابعاً: مع المنزل الذكي: الهاتف هو “المفتاح الرئيسي” للمنزل}

المنزل الذكي يبدو غامضاً، لكنه في جوهره ثلاثة أمور: التحكم عن بُعد عبر الهاتف، ووضع قواعد الأتمتة، والتحكم الصوتي.

دعني أحدثك عن بعض استخداماتي الفعلية، لتفهم مدى فائدته:

\begin{itemize}
  \item \textbf{المكيف}: في طريق العودة صيفاً، أشغل المكيف من الهاتف مسبقاً، فأجد المنزل بارداً عند الدخول. اشترِ مقبس مكيف مساعد بمئة يوان تقريباً لتحقيق ذلك، دون تغيير المكيف.
  \item \textbf{الإضاءة}: قبل النوم في السرير، كسلٌ للنهوض لإطفاء الضوء، انقر على الهاتف أو قل "يا Siri/Xiao Ai، أطفئ الضوء".
  \item \textbf{روبوت التنظيف}: قبل الخروج، انقر "ابدأ التنظيف" على الهاتف، وعند العودة يكون الأرض نظيفاً. ويمكن جدولة تنظيف يومي.
  \item \textbf{الكاميرا}: إذا كان لديك كبار السن أو أطفال أو حيوانات أليفة في المنزل، يمكنك رؤية ما يحدث في المنزل في أي وقت من الهاتف.
\end{itemize}

أي نظام بيئي تختار؟ خلاصة في جملة:

\begin{itemize}
  \item إذا كنت تستخدم iPhone أعطِ الأولوية للأجهزة التي تدعم Apple HomeKit (أغلى لكن حماية الخصوصية أفضل).
  \item إذا كانت أجهزتك كثيرة وتبحث عن القيمة مقابل السعر نظام Xiaomi Mi Home الأكثر اكتمالاً، فالإضاءة والمقابس وروبوتات التنظيف كلها متوفرة.
  \item إذا كنت تستخدم هاتف Huawei أعطِ الأولوية لـ HarmonyOS Connect، فالتجربة أفضل مع أجهزة Huawei.
\end{itemize}

\booksection{خامساً: سماعات ومكبرات صوت Bluetooth: حرية اللاسلك}

الانتقال من السماعات السلكية إلى سماعات Bluetooth كان أكثر تحسن ملموس في تجربة الهاتف بالنسبة لي. لن أخرج بعد الآن كرة من الأسلاك المتشابكة من الحقيبة، ولن أقلق من تعلق السلك بمقبض الباب.

\begin{itemize}
  \item \textbf{الإقران بسيط}: السماعات عند أول تشغيل تدخل عادةً وضع الإقران تلقائياً (وميض)، افتح Bluetooth على الهاتف لتجدها، وانقر للاتصال. بعد ذلك تتصل تلقائياً في كل مرة تخرج فيها، دون فعل أي شيء.
  \item \textbf{للعلامة نفسها مزايا إضافية}: AirPods مع iPhone تظهر نافذة منبثقة تعرض مستوى البطارية، ويمكن التبديل بسلاسة بين أجهزة Apple المتعددة. سماعات Huawei FreeBuds مع هواتف Huawei، وسماعات Xiaomi مع هواتف Xiaomi، لها نوافذ منبثقة وتحسينات مماثلة.
  \item \textbf{مكبر صوت Bluetooth}: في التجمعات العائلية أو رحلات التخييم، يمكن لمكبر صوت Bluetooth مضاعفة الأجواء. بعد الاتصال، يصبح الهاتف منصة الـ DJ الخاصة بك.
  \item \textbf{انتبه للتأخير}: عند اللعب بسماعات Bluetooth، قد يكون هناك تأخير طفيف (تسمع صوت الطلقة بعد نصف ثانية). إذا كنت حساساً للتأخير، فاستخدم سماعات سلكية، أو اشترِ سماعات Bluetooth تدعم "وضع التأخير المنخفض".
\end{itemize}

\begin{tcbinfo}{تقنية UWB فائقة العرض — "إحساس مكاني" أدق من البلوتوث}
إلى جانب البلوتوث، تخفي الهواتف الجديدة شريحة UWB (Ultra-Wideband، النطاق العريض جداً). قدرتها هي\textbf{قياس المسافة بدقة وتحديد الموقع المكاني}— تستطيع حساب مدى بُعد جهاز آخر عنك وفي أي اتجاه، بدقة تصل لمستوى السنتيمتر. البلوتوث يخبرك فقط “في الجوار”، لكن UWB يخبرك “على بُعد 1.2 متر إلى الأمام اليسار”. تطبيقات شائعة:

\begin{itemize}
  \item \textbf{AirTag/بحث دقيق}: عند بحث iPhone عن AirTag يعرض سهم تنقل "على بُعد 7 أمتار، امشِ لليسار"، هذا بفضل UWB. البلوتوث العادي يقول فقط "في النطاق أو لا".
  \item \textbf{مفتاح سيارة رقمي}: iPhone والأجهزة الرائدة من Android المزودة بـ UWB يمكن استخدامها كمفتاح سيارة، فتح القفل تلقائياً عند الاقتراب وقفله عند الابتعاد، وتمييز موقع السائق عن الراكب، ومقاومة الكسر بالترحيل.
  \item \textbf{الملاحة الداخلية}: في المولات والمطارات عبر منارات UWB لتحقيق ملاحة "على مستوى الممر"، أدق بكثير من GPS.
  \item \textbf{الموديلات المدعومة}: iPhone 11 والأحدث جميعها مزودة بشريحة UWB (U1/U2)؛ في معسكر Android، Xiaomi وSamsung وvivo وOPPO بعض الأجهزة الرائدة مزودة بها.
\end{itemize}

\end{tcbinfo}

\begin{tcbinfo}{LE Audio وAuracast — الجيل الجديد من معايير الصوت عبر البلوتوث}
الكلاسيكي الذي استُخدم لعشرين عاماً (A2DP) يحل محله\textbf{LE Audio}(صوت منخفض الطاقة)، وهو معيار الصوت الجديد منذ Bluetooth 5.2:

\begin{itemize}
  \item \textbf{أكثر توفيراً للطاقة وجودة صوت أفضل}: يستخدم ترميز LC3 الجديد، جودة صوت أفضل من SBC القديم واستهلاك طاقة أقل، وعمر بطارية السماعة أطول بوضوح.
  \item \textbf{صوت ثنائي القناة}: كل سماعة (يسار/يمين) تنقل بشكل مستقل، التبديل لوضع سماعة واحدة لا يقطع الاتصال.
  \item \textbf{بث Auracast الصوتي}: الوظيفة المميزة لـ LE Audio. هاتف/تلفزيون واحد يمكنه البث في آن واحد إلى\textbf{أعداد لا نهائية من}سماعات الرأس — في صالات انتظار المطارات والمتاحف والصالات الرياضية والفصول الدراسية يمكن لمصدر إرسال واحد أن يخدم عدة أشخاص بسماعاتهم الخاصة يستمعون لنفس المحتوى.
  \item \textbf{الوضع الحالي}: سماعات الرأس والهواتف الرائدة الجديدة 2024-2026 تدعم LE Audio بشكل عام، الأجهزة القديمة تحتاج ترقية عتاد ولا يمكن فتحها عبر البرامج. عند شراء سماعة جديدة انتبه لعلامة دعم LE Audio/Auracast.
\end{itemize}

\end{tcbinfo}

\booksection{سادساً: مع نظام السيارة: شريك آمن أثناء القيادة}

استخدام الهاتف أثناء القيادة غير قانوني، وقاتل. لكن أحياناً تحتاج فعلاً إلى الملاحة أو الرد على المكالمات أو تشغيل الموسيقى. هنا يأتي دور CarPlay (Apple) وAndroid Auto (أندرويد).

\begin{itemize}
  \item \textbf{كيف تتصل}: إذا كانت سيارتك تدعم ذلك (معظم السيارات الحديثة تدعم)، استخدم كابل بيانات لتوصيل الهاتف بمنفذ USB في السيارة، فتتحول شاشة السيارة تلقائياً إلى واجهة CarPlay أو Android Auto. بعض السيارات الجديدة تدعم الاتصال اللاسلكي، وهو أسهل.
  \item \textbf{ماذا يمكنك أن تفعل}: الملاحة (أفضل بكثير من ملاحة السيارة)، وإجراء المكالمات (الرد من عجلة القيادة)، والاستماع للموسيقى (تطبيقات الموسيقى على الهاتف تُعرض مباشرة على السيارة)، وإرسال رسائل WeChat صوتياً (تقرأ Siri أو Xiao Ai الرسائل وترد عليها).
  \item \textbf{السيارات القديمة أيضاً}: إذا كانت سيارتك قديمة ولا تدعم، اشترِ بشيء يوان شاشة CarPlay خارجية، الصقها على لوحة العدادات، والتجربة مشابهة.
\end{itemize}

\booksection{سابعاً: الأمان والخصوصية في الربط}

كلما زادت الأجهزة المربوطة، زاد سطح خطر تسريب الخصوصية. بعض النقاط التي يجب الانتباه لها:

\begin{itemize}
  \item قبل العرض على الشاشة في الأماكن العامة، تحقق من شريط إشعارات الهاتف لعدم وجود رسائل خاصة، حتى لا يرى الجميع ما لا يجب رؤيته.
  \item قبل إعطاء الجهاز لشخص آخر أو بيعه، أزل الإقران — اذهب إلى إعدادات الهاتف، ابحث عن ذلك الجهاز، واختر "إزالة" أو "إلغاء الإقران".
  \item كاميرات المنزل الذكي، اشترِ ذات علامة تجارية ومشفّرة، لا تشترِ الرخيصة بعشرات اليوانات. سمعنا كثيراً أخباراً عن كاميرات تم بثّها من قبل القراصنة.
\end{itemize}

\begin{tcbpractice}{جرّب: اعرض شاشة هاتفك على التلفاز، وأعد مشهداً ذكياً}
\begin{enumerate}
  \item تأكد من اتصال الهاتف والتلفاز بنفس شبكة Wi-Fi.
  \item ابحث عن ميزة "Screen Mirroring" أو "AirPlay" على التلفاز وفعّلها.
  \item على الهاتف: iPhone انقر "Screen Mirroring" من مركز التحكم؛ أندرويد اسحب للأسفل وابحث عن "العرض اللاسلكي".
  \item اختر صورتك المفضلة واعرضها على التلفاز ليراها أفراد العائلة.
  \item إذا كان لديك أجهزة منزل ذكي، حاول إنشاء "وضع مغادرة المنزل" في تطبيق الهاتف: إطفاء الأضواء وإطفاء المكيف وتشغيل روبوت التنظيف بنقرة واحدة.
\end{enumerate}

\end{tcbpractice}

\begin{tcbinfo}{حالة 1: تحسين عمر البطارية عملياً}
\begin{itemize}
  \item \textbf{السيناريو}: مستخدم Xiaomi 14 عمر بطارية ضعيف، ينفد في نصف يوم.
  \item \textbf{الخطوات}: (1) الإعدادات البطارية عرض ترتيب الاستهلاك (2) تحديد التطبيقات الخلفية غير الطبيعية (3) إغلاق 5G والتبديل لـ 4G (4) تقييد الخلفية (5) تفعيل وضع توفير الطاقة (6) نزع الغطاء أثناء الشحن لتبريد الجهاز.
  \item \textbf{النتيجة المتوقعة}: إطالة عمر البطارية بنسبة 30-50\%، يمكن قراءة بيانات صحة البطارية.
\end{itemize}

\end{tcbinfo}
